\section{Background and Related Work}
\label{sec:background}

\subsection{Preservation as a transformation requirement}
ATAG 2.0's transformation guidance treats restructuring and recoding as accessibility-relevant operations: when the destination technology provides an equivalent mechanism, accessibility information should be preserved or the author should be warned or prompted to check it \cite{atag20}. Revised Section 508 identifies preservation of information provided for accessibility in format conversion, and EN 301 549 contains a related transformation requirement \cite{section508,etsi301549}. These sources motivate the question but do not predict the feature-specific outcomes of the pipelines evaluated here.

\subsection{Office-document conversion}
Roig and Ribera studied accessibility consequences of office-document conversion to EPUB \cite{roig2015}. Their work is a direct precedent for treating conversion as an accessibility question rather than a purely syntactic operation. A11yPreserve differs in destination format, source instrumentation, and decision boundary: it uses frozen DOCX contracts, raw OOXML checks, named DOCX-to-PDF pipelines, PDF structure extraction, provenance records, and an explicit unresolved category. The distinction is one of scope and method, not a claim that prior work ignored preservation.

\subsection{Oracle construction and transformation testing}
The oracle problem is a long-standing software-testing concern: executing a test does not guarantee that the output can be judged correct \cite{barr2015oracle}. A11yPreserve uses a specification-derived feature oracle. The source contract states the intended property, while the destination comparator must determine whether a PDF representation is accepted as equivalent. Differential testing supplies the source/destination pairing, but the comparison is not a majority vote across converters. Metamorphic testing is a related strategy for constructing expected relations \cite{segura2016metamorphic}; it is not the primary oracle here because the source contracts provide explicit expected properties. Format-conversion scholarship likewise emphasizes that successful conversion requires a declared notion of what must be preserved \cite{sperberg2011conversion}.

\subsection{Accessible destination generation and evaluation}
SciA11y converts scientific PDFs to accessible HTML and focuses on navigation-oriented destination generation \cite{wang2021scia11y}. iTagPDF uses source cues and PDF analysis to improve tagging, reading order, alternative text, and formula structure \cite{mowar2026}. Kumar et al. benchmark automated and LLM-based PDF accessibility evaluation across criteria including alternative text, reading order, tables, and links \cite{kumar2025}. DAISY's benchmark compares PDF-to-accessible-document services using shared criteria and consensus scoring \cite{daisy2026}. DocAccessible reports source-grounded PDF-to-HTML benchmarking and practical Google Docs export guidance \cite{docaccessible2026,docaccessible-google}. These efforts provide important context, but they answer different questions or use different directions of conversion.

A11yPreserve's boundary is therefore modest: it does not invent accessibility-conversion benchmarking, replace destination-only validators, or claim a complete accessibility checker. It supplies a controlled source-to-destination preservation protocol whose evidence boundary is explicit and whose uncertainty is retained.
